\section{Gain Design}

The gain schedule was computed at Mach \(0.9\), \(10~\text{m}\) altitude, and four mass/CG points. The internal gain is the same for the attitude and acceleration autopilots because both use the same \(q\)-feedback loop.

\begin{table}[H]
\centering
\caption{Computed gain schedule.}
\begin{tabular}{rrrrrr}
\toprule
\(t\) (s) & \(X_{CG}\) (m) & \(K_q\) attitude & \(K_\theta\) & \(K_q\) accel. & \(K_a\) \\
\midrule
0.000 & 2.8325 & 0.6000 & 9.3333 & 0.6000 & 0.0250 \\
1.725 & 2.8934 & 0.6667 & 10.6667 & 0.6667 & 0.0250 \\
3.450 & 2.9647 & 0.7333 & 12.0000 & 0.7333 & 0.0250 \\
115.350 & 2.6012 & 0.3333 & 8.0000 & 0.3333 & 0.0250 \\
\bottomrule
\end{tabular}
\end{table}

\begin{figure}[H]
\centering
\includegraphics[width=0.95\textwidth]{figures/gain_curves_vs_cg.png}
\caption{Internal and external gains as a function of CG position.}
\end{figure}

\begin{figure}[H]
\centering
\includegraphics[width=0.9\textwidth]{figures/gain_schedule_vs_time.png}
\caption{Gain schedule along the representative burn points.}
\end{figure}

\subsection{Root-Locus Plots}

The following root-locus plots show the nominal design point, taken during the sustainer phase. The internal loop is common to both autopilots.

\begin{figure}[H]
\centering
\includegraphics[width=0.78\textwidth]{figures/root_locus_inner_attitude.png}
\caption{Root locus of the internal attitude loop.}
\end{figure}

\begin{figure}[H]
\centering
\includegraphics[width=0.78\textwidth]{figures/root_locus_outer_attitude.png}
\caption{Root locus of the external attitude loop.}
\end{figure}

\begin{figure}[H]
\centering
\includegraphics[width=0.78\textwidth]{figures/root_locus_inner_acceleration.png}
\caption{Root locus of the internal loop used by the acceleration autopilot.}
\end{figure}

\begin{figure}[H]
\centering
\includegraphics[width=0.78\textwidth]{figures/root_locus_outer_acceleration.png}
\caption{Root locus of the external acceleration loop.}
\end{figure}
